\section{How large could sampling bias be?}
\label{sec:scope}

Before developing machinery, we establish the scale of the problem. This section derives a
worst-case bound on sampling bias in a regional ocean heat content (OHC) average that depends
on a single interpretable quantity---how unevenly floats sample the region---and requires no
model for the ocean beyond quantities computable from the profiles themselves. The bound is
deliberately crude. Its purpose is to determine whether the question is live, and to identify
what must be measured for the remainder of the paper to have content.

\subsection{Two kinds of bias}
\label{sec:two-kinds}

It is useful to separate two mechanisms that are often treated together.

\emph{Coverage bias} arises because floats are unevenly distributed in space and time: more
profiles in summer than winter, more in the subtropics than at high latitudes, more in 2019
than in 2004. Since the ocean's mean state varies with position and season, weighting
positions incorrectly biases the average. Crucially, coverage bias is visible in the sampling
geometry alone. Given a map of profile locations and a climatology, it can be computed and
removed.

\emph{Preferential sampling bias} arises when the positions themselves depend on the field
being measured \citep{diggle2010}. It persists even when the distribution of positions is
perfectly uniform. Consider a region sampled at uniform density in which one float is
entrained into a warm-core anticyclone and remains for four months, returning twelve profiles,
while a neighbouring float in a cold cyclone is advected out of the region within three weeks,
returning two. Profile density is uniform; the profiles nonetheless over-represent warm water.
No amount of reweighting by position corrects this, because the imbalance is invisible in the
positions.

The distinction matters for Argo specifically because floats are advected by ocean currents,
and currents are geostrophically related, through thermal wind, to the density structure that
determines heat content. The instrument is transported by the field it measures. Whether this
coupling produces a material bias is an empirical question; this section establishes how
material it could be.

\subsection{Setup}
\label{sec:setup}

Let $\mathcal{R} = D \times \mathcal{T}$ denote a spatio-temporal analysis region, with $D$ a
horizontal domain restricted to bathymetry deeper than $2000$\,m and $\mathcal{T}$ an
averaging window. Let $P$ be the normalised uniform (area $\times$ time) measure on
$\mathcal{R}$, and let $Y : \mathcal{R} \to \mathbb{R}$ denote $0$--$2000$\,m OHC in
J\,m$^{-2}$. The target is
\begin{equation}
  \bar{Y} \;=\; \int_{\mathcal{R}} Y \, dP .
\end{equation}
Argo delivers profiles at locations $z_1, \dots, z_n \in \mathcal{R}$, with empirical sampling
measure $Q_n = n^{-1}\sum_{j=1}^{n} \delta_{z_j}$ and estimator
$\hat{Y} = \int Y \, dQ_n = n^{-1}\sum_j Y(z_j)$. The error is
\begin{equation}
  \hat{Y} - \bar{Y} \;=\; \int_{\mathcal{R}} Y \, d(Q_n - P).
\end{equation}

Following standard practice \citep{roemmich2009, cressie2011}, decompose
\begin{equation}
  \label{eq:decomp}
  Y \;=\; m \;+\; A \;+\; \varepsilon ,
\end{equation}
where $m$ is a climatological mean field taken as known from an independent long record, $A$
is a resolved-scale transient anomaly lying in a finite-dimensional space
$V_L \subset L^2(P)$ containing the constants, and $\varepsilon$ is the unresolved remainder.
We take the split to be a definition rather than an approximation: with $\Pi_L$ the
$L^2(P)$-orthogonal projection onto $V_L$,
\begin{equation}
  A := \Pi_L (Y - m), \qquad \varepsilon := (I - \Pi_L)(Y-m),
\end{equation}
so that \eqref{eq:decomp} holds exactly. Because $\mathbf{1} \in V_L$,
\begin{equation}
  \label{eq:zeromean}
  \int_{\mathcal{R}} \varepsilon \, dP \;=\; \langle \varepsilon, \mathbf{1}\rangle_{L^2(P)} \;=\; 0 .
\end{equation}

The error then separates as
\begin{equation}
  \label{eq:split}
  \hat{Y} - \bar{Y}
  \;=\; \underbrace{\int (m + A)\, d(Q_n - P)}_{\text{coverage}}
  \;+\; \underbrace{\int \varepsilon \, dQ_n}_{\text{preferential}},
\end{equation}
where the second term has been simplified using \eqref{eq:zeromean}. The coverage term
involves only structure that is either known ($m$) or estimable at resolved scales ($A$), and
is computable from the profile positions; it is addressed in Section~\ref{sec:coverage-later}.
The preferential term is the subject of this section.

\begin{remark}
\label{rem:sigma-definition}
Throughout, $\varepsilon$ denotes the residual field after removal of \emph{both} the
climatology and the resolved-scale anomaly, and $\sigma$ denotes its standard deviation under
$P$. This is substantially smaller than the standard deviation of raw OHC, and smaller than
that of the raw anomaly $Y - m$. Substituting either of the latter would inflate every bound
below and would conflate the two bias mechanisms.
\end{remark}

\subsection{Why the preferential term cannot be estimated from the profiles}
\label{sec:not-estimable}

The preferential term $\int \varepsilon\, dQ_n = n^{-1}\sum_j \varepsilon(z_j)$ is the mean of
the residual field \emph{at the sampled locations}. It would be zero if the sampled locations
were a fair draw with respect to $\varepsilon$, and is nonzero exactly to the extent that they
are not.

It cannot be estimated directly. The residuals available to us, $\hat\varepsilon_j$, are
obtained by fitting $m + A$ to the observed profiles; when $\mathbf{1}$ is in the fitted span,
the fitted residuals satisfy $\sum_j \hat\varepsilon_j = 0$ identically, by construction. More
fundamentally, $\varepsilon$ is observed only where floats went. If floats systematically
avoided cold water, no feature of the sample would reveal the temperature of the water they
avoided.

The preferential term must therefore be \emph{bounded under an assumption}, not estimated. The
remainder of this section makes that assumption explicit, derives the sharpest bound it
supports, and shows how the assumption itself can be checked against independent observations.

\subsection{The sampling-ratio assumption}
\label{sec:lambda}

Let $Q$ denote the limiting sampling measure on $\mathcal{R}$ (the intensity of the profile
point process, normalised to unit mass), so that $Q_n \to Q$ as the array is sampled more
densely. We assume that no part of the region is grossly over- or under-sampled relative to
its area-time share.

\begin{assumption}[Bounded sampling ratio]
\label{ass:lambda}
There exists $\Lambda \ge 1$ such that $Q$ is absolutely continuous with respect to $P$ with
\begin{equation}
  \Lambda^{-1} \;\le\; \frac{dQ}{dP}(z) \;\le\; \Lambda
  \qquad \text{for $P$-almost every } z \in \mathcal{R}.
\end{equation}
\end{assumption}

Two features deserve emphasis.

First, the assumption is \emph{two-sided}, and the lower bound is what makes the problem
tractable. A one-sided assumption permits floats to miss some water types entirely, in which
case nothing about those water types is recoverable and no finite bound exists. The lower
bound asserts that every water type is present in the sample---floats do enter cold cyclones,
they simply leave sooner---so that the truth is a \emph{reweighting} of the observed sample
rather than a superset of it. This is the substantive content of the assumption and the point
at which it is most vulnerable (Section~\ref{sec:limitations}).

Second, Assumption~\ref{ass:lambda} is stated over the region \emph{after} the coverage term
has been accounted for; that is, $\Lambda$ quantifies residual non-uniformity not already
explained by the resolved structure. In practice this is enforced by measuring $\Lambda$ over
a partition constructed from mesoscale features (Section~\ref{sec:measuring}) rather than from
large-scale geography.

\subsection{Worst-case bias}
\label{sec:bound}

Let $F$ denote the distribution of the value $\varepsilon(z)$ when $z \sim P$, and $G$ the
corresponding distribution when $z \sim Q$. Assumption~\ref{ass:lambda} transfers to these
push-forward laws: for any Borel set $B$, $Q(\varepsilon^{-1}B) \in [\Lambda^{-1}P(\varepsilon^{-1}B),\,
\Lambda P(\varepsilon^{-1}B)]$, hence $dG/dF \in [\Lambda^{-1}, \Lambda]$. By
\eqref{eq:zeromean}, $\mathbb{E}_F[\varepsilon] = 0$, and the preferential bias is
$\mathbb{E}_G[\varepsilon]$.

Define the \emph{upper partial expectation} of $F$ at level $p \in (0,1)$,
\begin{equation}
  \Psi(p) \;:=\; \int_{\{e \,>\, F^{-1}(1-p)\}} e \; dF(e)
  \;=\; p \cdot \mathbb{E}_F\!\left[\varepsilon \,\middle|\, \varepsilon > F^{-1}(1-p)\right],
\end{equation}
the contribution to the mean made by the largest $p$ fraction of the field.

\begin{proposition}[Sharp worst-case preferential bias]
\label{prop:bound}
Under Assumption~\ref{ass:lambda},
\begin{equation}
  \label{eq:bound}
  \left| \int \varepsilon \, dQ \right|
  \;\le\; B(\Lambda) \;:=\; \left(\Lambda - \Lambda^{-1}\right)\,
  \Psi\!\left(\frac{1}{\Lambda + 1}\right),
\end{equation}
and the bound is attained.
\end{proposition}

\begin{proof}
Write $r = dG/dF$. The problem is
\[
  \max_{r} \int e \, r(e) \, dF(e)
  \quad \text{s.t.} \quad \int r \, dF = 1, \quad \Lambda^{-1} \le r \le \Lambda,
\]
a linear program in $r$. The objective is increasing in $r(e)$ where $e > 0$ and decreasing
where $e < 0$, and the single equality constraint is a total-mass budget; by the bathtub
principle the optimum is bang-bang, $r = \Lambda$ on $\{e > q\}$ and $r = \Lambda^{-1}$ on
$\{e < q\}$ for a threshold $q$. Writing $p = F(\{e > q\})$, the mass constraint gives
\[
  \Lambda p + \Lambda^{-1}(1 - p) = 1
  \quad \Longrightarrow \quad
  p = \frac{1 - \Lambda^{-1}}{\Lambda - \Lambda^{-1}} = \frac{1}{\Lambda + 1}.
\]
Since $\mathbb{E}_F[\varepsilon] = 0$ we have $\int_{\{e < q\}} e \, dF = -\Psi(p)$, so the
optimal value is $\Lambda \Psi(p) - \Lambda^{-1}\Psi(p) = (\Lambda - \Lambda^{-1})\Psi(p)$.
The lower bound follows by symmetry (interchanging the roles of the two levels), giving
$|\mathbb{E}_G[\varepsilon]| \le B(\Lambda)$.
\end{proof}

Note $B(1) = 0$, as required: if sampling is exactly proportional to area-time, there is no
preferential bias.

\subsection{A distribution-free first-order form}
\label{sec:first-order}

Expression \eqref{eq:bound} requires the tail structure of $F$. For $\Lambda$ near unity it
simplifies considerably. As $\Lambda \to 1$ we have $\Lambda - \Lambda^{-1} = 2(\Lambda - 1) +
O((\Lambda-1)^2)$ and $p \to 1/2$, and for any distribution symmetric about zero,
$\Psi(1/2) = \tfrac12 \mathbb{E}_F|\varepsilon|$. Hence:

\begin{corollary}
\label{cor:firstorder}
To first order in $\Lambda - 1$, and for $\varepsilon$ symmetric under $P$,
\begin{equation}
  \label{eq:firstorder}
  \boxed{\;\left| \int \varepsilon \, dQ \right| \;\lesssim\; (\Lambda - 1)\;
  \mathbb{E}_P\left|\varepsilon\right| \;}
\end{equation}
\end{corollary}

This is the interpretable form of the result: \emph{the worst-case preferential bias is
approximately the fractional sampling imbalance times the mean absolute residual}. Both
factors are concrete. The second is computed directly from the profiles by averaging
$|\hat\varepsilon_j|$; it requires no distributional assumption and is robust to the skewness
and heavy tails typical of mesoscale variability. The first is the quantity to be measured in
Section~\ref{sec:measuring}.

For a Gaussian residual field, $\mathbb{E}|\varepsilon| = \sigma\sqrt{2/\pi} \approx
0.798\,\sigma$, so \eqref{eq:firstorder} reads $B \approx 0.8\,(\Lambda-1)\,\sigma$.

\subsection{Magnitudes}
\label{sec:magnitudes}

Table~\ref{tab:lambda} evaluates \eqref{eq:bound} exactly for a Gaussian residual field. The
Gaussian case is used only for orientation; in application, $\Psi$ is evaluated empirically
from the sorted residuals.

\begin{table}[t]
\centering
\begin{tabular}{lccc}
\toprule
$\Lambda$ & $B(\Lambda)/\sigma$ & \multicolumn{2}{c}{$B(\Lambda)$, J\,m$^{-2}$} \\
\cmidrule(lr){3-4}
 & & $\sigma = 10^{9}$ & $\sigma = 3\times10^{8}$ \\
\midrule
$1.02$ & $0.016$ & $1.6\times10^{7}$ & $4.7\times10^{6}$ \\
$1.05$ & $0.039$ & $3.9\times10^{7}$ & $1.2\times10^{7}$ \\
$1.10$ & $0.076$ & $7.6\times10^{7}$ & $2.3\times10^{7}$ \\
$1.15$ & $0.111$ & $1.1\times10^{8}$ & $3.3\times10^{7}$ \\
$1.20$ & $0.145$ & $1.5\times10^{8}$ & $4.4\times10^{7}$ \\
$1.25$ & $0.178$ & $1.8\times10^{8}$ & $5.3\times10^{7}$ \\
$1.50$ & $0.322$ & $3.2\times10^{8}$ & $9.7\times10^{7}$ \\
$2.00$ & $0.545$ & $5.5\times10^{8}$ & $1.6\times10^{8}$ \\
\bottomrule
\end{tabular}
\caption{Worst-case preferential sampling bias from Proposition~\ref{prop:bound}, evaluated
for a Gaussian residual field. The two right-hand columns correspond to a mesoscale-energetic
region and a quiescent region respectively. For reference, a $0.01\,^\circ$C mean anomaly over
$2000$\,m is $\approx 8.2\times10^{7}$\,J\,m$^{-2}$ (taking $\rho c_p = 4.1\times10^{6}$
J\,m$^{-3}$K$^{-1}$), and multi-year regional OHC anomaly signals in the North Atlantic are of
order $10^{8}$\,J\,m$^{-2}$.}
\label{tab:lambda}
\end{table}

Two conclusions follow, and they set the agenda for the paper.

\emph{The question is live.} In a mesoscale-energetic region, a sampling imbalance of $15\%$
suffices to generate a bias comparable to the signal. Inverting \eqref{eq:bound}, the critical
value at which the bound equals a $10^{8}$\,J\,m$^{-2}$ signal is $\Lambda^\dagger \approx
1.13$ for $\sigma = 10^{9}$\,J\,m$^{-2}$. Imbalances of this size are not obviously
implausible: the Lagrangian coherence of mesoscale eddies is well documented, and anticyclones
are both warm-core and comparatively effective at retaining floats.

\emph{The answer is regionally specific.} For $\sigma = 3\times10^{8}$\,J\,m$^{-2}$, the
critical value is $\Lambda^\dagger \approx 1.5$, a far weaker requirement. Because $B$ scales
linearly in $\sigma$ while $\sigma$ varies by a factor of several between quiescent and
frontal regions, the same sampling behaviour is benign in one setting and disqualifying in
another. This motivates the paired-region design adopted in
Section~\ref{sec:design-later}.

\begin{remark}[Averaging does not remove this term]
\label{rem:no-averaging}
It is tempting to argue that a multi-year average washes out eddy-scale sampling effects. It
would, if the imbalance changed sign at random: bias from $N$ independently signed encounters
would decline as $N^{-1/2}$. The concern here is that the leading mechanism is
\emph{sign-definite}---warm anticyclones preferentially retain floats, in the same direction,
in every season. A systematic component of $\Lambda$ does not decline with the length of the
averaging window or with the size of the array. This is the essential difference between the
preferential term and the noise term of Section~\ref{sec:noise-later}, which does average
down.
\end{remark}

\subsection{Measuring $\Lambda$}
\label{sec:measuring}

Assumption~\ref{ass:lambda} is only useful if $\Lambda$ can be constrained observationally.
Satellite altimetry provides the required independent view: mesoscale eddy identification and
tracking atlases derived from sea surface height \citep{chelton2011} give, for each day, the
polarity and boundary of every detected eddy, using measurements entirely independent of the
Argo profiles.

Partition $\mathcal{R}$ into classes $\kappa \in \{\text{anticyclone interior},\
\text{cyclone interior},\ \text{exterior}\}$ and compute
\begin{equation}
  f_\kappa \;=\; P(\text{class }\kappa),
  \qquad
  \hat{f}_\kappa \;=\; Q_n(\text{class }\kappa),
\end{equation}
the areal--temporal fraction occupied by each class and the fraction of profiles falling
within it. Under exogenous sampling these agree up to sampling error. The estimator
\begin{equation}
  \label{eq:lambdahat}
  \hat{\Lambda} \;=\; \max_{\kappa} \; \max\!\left\{ \frac{\hat{f}_\kappa}{f_\kappa},\;
  \frac{f_\kappa}{\hat{f}_\kappa} \right\}
\end{equation}
then measures the sampling imbalance at the resolution of the partition. If, for example,
anticyclone interiors occupy $20\%$ of the region and contain $22\%$ of profiles,
$\hat\Lambda = 1.10$.

Two qualifications are essential.

\emph{$\hat\Lambda$ is a lower bound on $\Lambda$.} Preferential structure operating within
classes---selection by eddy amplitude, by position relative to the eddy core, or at
submesoscale---is invisible to a coarse partition. The remedy is to refine the partition
(stratifying additionally by eddy amplitude quantile and by normalised radius) and report
$\hat\Lambda$ as a function of refinement. Saturation indicates that the relevant structure
has been captured; continued growth indicates that it has not.

\emph{$\hat\Lambda$ carries substantial sampling error.} Because consecutive profiles from one
float are separated by $30$--$80$\,km---well inside a mesoscale feature---they are close to
replicates, and the effective sample size is far below the profile count. With $n_{\text{eff}}
\sim 300$ and $f_\kappa \approx 0.2$, the standard error on $\hat f_\kappa$ is
$\approx 0.023$, giving a standard error on the ratio of roughly $0.12$. This is comparable to
$\Lambda^\dagger - 1$ itself. Constraining $\Lambda$ tightly enough to be decisive therefore
requires pooling across the full record and the cleanest available partition, and the
resulting uncertainty must be propagated into \eqref{eq:bound} rather than reported alongside
it.

\subsection{A tighter bound under class-restricted preferentiality}
\label{sec:refinement}

Proposition~\ref{prop:bound} is deliberately agnostic: it permits the sampler to concentrate
on whatever part of the field is warmest, which is why the worst case involves a tail mean. If
one is willing to assume that preferential sampling operates \emph{only} through the observed
partition---that is, that sampling is exogenous conditional on eddy class---the same linear
program over the three-cell partition yields
\begin{equation}
  \label{eq:stratified}
  \left| \int \varepsilon \, dQ \right| \;\le\; \sum_{\kappa}
  \left| \hat{f}_\kappa - f_\kappa \right| \cdot \left| \bar\varepsilon_\kappa \right| ,
\end{equation}
where $\bar\varepsilon_\kappa$ is the mean residual within class $\kappa$, estimable from the
profiles falling in that class. Because \eqref{eq:stratified} involves class \emph{means}
rather than tail means, it is typically a factor of three to five tighter than
\eqref{eq:bound}.

The two bounds are complementary and should both be reported: \eqref{eq:bound} is
assumption-light and pessimistic, \eqref{eq:stratified} is sharper but conditional on the
partition capturing the mechanism. Their ratio is itself informative about how much is being
bought by the additional assumption. Expression \eqref{eq:stratified} also suggests an
immediate remedy---post-stratifying the profiles so that $\hat f_\kappa = f_\kappa$ removes
the class-driven component of the bias by construction.

\subsection{Limitations}
\label{sec:limitations}

\emph{Regions of zero coverage.} If some part of $\mathcal{R}$ is never sampled---beneath
permanent or seasonal ice, or at depths shallower than $2000$\,m where the target is
undefined---then $dQ/dP = 0$ there, Assumption~\ref{ass:lambda} fails, and no finite bound is
available. Such regions must be excised from $\mathcal{R}$ and discussed separately rather
than folded into an average. This is the assumption's principal failure mode, and it is severe
in the seasonally ice-covered Southern Ocean.

\emph{The empirical distribution is $G$, not $F$.} Evaluating $\Psi$ from the observed
residuals estimates the tail structure of the sampled distribution rather than the true one.
Since the two differ by a density ratio in $[\Lambda^{-1},\Lambda]$, the induced error in
\eqref{eq:bound} is $O((\Lambda-1)^2)$ and is negligible in the regime of interest; a fully
conservative treatment substitutes $\sqrt{\Lambda}\,\hat\sigma$ for $\sigma$, inflating the
bound by $\sqrt{\Lambda}$ (approximately $5\%$ at $\Lambda = 1.1$).

\emph{Dependence on the resolved/unresolved split.} The magnitude of $\sigma$, and hence of
the bound, depends on the choice of $V_L$. A larger resolved space transfers variance from
$\varepsilon$ into $A$, shrinking the preferential bound while enlarging the coverage term. The
split must therefore be fixed by an independent criterion and its sensitivity reported;
Section~\ref{sec:truncation-later} does this by requiring that the residual field have
collapsed to short-range correlation.

\emph{The bound is worst-case, not an estimate.} Proposition~\ref{prop:bound} states what
sampling imbalance \emph{could} do, not what it does. A small bound is therefore informative
(it rules the mechanism out) while a large one is not (it establishes only that the mechanism
is not ruled out). Establishing the actual magnitude requires the further work of
Sections~\ref{sec:osse-later} and \ref{sec:budget-later}.

\subsection{Implications for what follows}
\label{sec:roadmap}

Table~\ref{tab:lambda} and the critical values $\Lambda^\dagger \approx 1.13$ and $1.5$ define
the target precision for the rest of the paper. Three requirements follow.

First, $\sigma$---the standard deviation of the residual after removal of climatology and
resolved anomaly---must be estimated, which requires specifying $V_L$. Second, $\Lambda$ must
be constrained to a precision finer than $\Lambda^\dagger - 1$, which given the effective
sample size is demanding and motivates careful treatment of the Lagrangian replication
structure of the profile record. Third, because the bound is worst-case, an independent
estimate of the actual bias is needed, which we obtain by advecting synthetic floats through
an eddy-resolving circulation model and comparing against the model's known heat content.

We take these in turn, and report throughout the pair of contrasting regions motivated by the
$\sigma$-dependence of Table~\ref{tab:lambda}.
